\section{结论}

Google File System展现了在普通硬件上支撑大规模数据处理工作负载所必需的特性。
尽管某些设计决策是针对我们独特环境作出的，但其中许多决策也适用于规模和成本敏感
度相近的数据处理任务。


我们从当前及预期的应用工作负载和技术环境出发，重新审视了传统文件系统的假设。
这些观察使我们在设计空间中选择了截然不同的方案。我们将组件故障视为常态而非例外，
针对主要以追加方式写入(可能是并发追加)、随后读取(通常为顺序读取)的巨大文件进行优化，
并通过扩展和放宽标准文件系统接口来改进整个系统。

我们的系统通过持续监控、复制关键数据以及快速自动恢复来提供容错能力。
chunk复制使我们能够容忍chunkserver故障。这些故障的频率促使我们设计了一种新的
在线修复机制，它会定期且透明地修复损坏，并尽快补偿丢失的replica。此外，我们使用校验
和检测磁盘或IDE子系统层面的数据损坏；考虑到系统中的磁盘数量，这类问题已经变得十分常见。

我们的设计能够向执行各种任务的大量并发读取者和写入者提供高聚合吞吐量。
我们通过将经由master的文件系统控制，与直接在chunkserver和客户端之间进行
的数据传输分离，实现这一目标。大chunk大小和chunk lease减少了master对常见
操作的参与；其中chunk lease将数据mutation中的权限委托给primary replica。
这使得我们能够采用简单的中心化master，同时避免其成为瓶颈。我们相信，对网络栈的
改进将消除当前单个客户端写入吞吐量的限制。

GFS已成功满足我们的存储需求，并在Google内部广泛用作研究与开发以及生产数据处理的
存储平台。它是一项重要工具，使我们能够持续创新，并解决整个Web规模的问题。
